% SOURCE_AUTHORITY: sga5_fr_workpass.tex, lines 14405--14415 (LF-normalized unit).
% SOURCE_SHA256: 791F4EFFC5E02832D5D77ED03518C8156D6F07E4C8238B03545DB93D883FBB28
Suponha-se agora que $X$ seja étale sobre $S$; o mesmo ocorre com $X^{(p)}$ por mudança de base e, por conseguinte, o morfismo $\Fr_{X/S}:X\to X^{(p)}$ é étale. Reciprocamente, suponha-se que $\Fr_{X/S}$ seja étale; ele é não ramificado, logo $X$ é não ramificado sobre $S$. Em outras palavras, o morfismo estrutural $g:X\to S$ é, localmente sobre $X$ para a topologia étale, uma imersão fechada (SGA I cor. 7.8). Pode-se, portanto, supor que $g$ seja uma imersão fechada de apresentação finita, e trata-se de ver que, se $\Fr_{X/S}$ for um isomorfismo, $g$ será uma imersão aberta. Ora, $X$ é definido por um ideal $\JJ$ de tipo finito,
\[
X=(|X|,\OO_S/\JJ),
\]
e, segundo b),
\[
X^{(p)}=(|X|,\OO_S/\JJ^p),
\qquad
\Fr_{X/S}=(\id_{|X|},\OO_S/\JJ^p\to\OO_S/\JJ).
\]
Dizer que $\Fr_{X/S}$ é um isomorfismo equivale, portanto, a dizer que $\JJ=\JJ^p$; como $\JJ$ é de tipo finito, deduz-se pelo lema de Nakayama que $\JJ_s=0$, ou $\JJ_s=\OO_{S,s}$, para todo ponto $s\in S$, de onde resulta a conclusão.
